%% Cover letter for IEEE Access submission
%% Compile with pdflatex
\documentclass[11pt, a4paper]{letter}
\usepackage{geometry}
\geometry{top=2.5cm, bottom=2.5cm, left=3cm, right=3cm}
\usepackage{url}
\usepackage{microtype}
\usepackage{parskip}

\signature{Florin Olariu\\
  Faculty of Computer Science\\
  Alexandru Ioan Cuza University of Ia\c{s}i, Romania\\
  \texttt{folariu@suvoda.com}\\[4pt]
  \small(on behalf of both authors)}

\address{Florin Olariu\\
  Faculty of Computer Science\\
  Alexandru Ioan Cuza University of Ia\c{s}i\\
  Bd.\ Carol I, Nr.\ 11, 700506 Ia\c{s}i\\
  Romania}

\date{\today}

\begin{document}

\begin{letter}{%
  Editor-in-Chief\\
  \textit{IEEE Access}
}

\opening{Dear Editor,}

We submit for consideration in \textit{IEEE Access} our manuscript,

\begin{center}
  \textbf{``A Skill-Mediated LLM Workflow for Migrating a .NET Monolith
  to Microservices: A Controlled Case Study''}
\end{center}

\textbf{What the paper reports.}
We present a controlled case study of a skill-mediated LLM workflow
applied end-to-end to a single .NET~10 monolith.  The workflow
composes ten Claude~Code skill scripts into a complete sequence:
requirements extraction, architecture analysis, microservices
decomposition, code generation, unit and integration testing,
automated review, automated security checks, automated remediation,
and gated publish.  The case study uses four publicly-available
artefacts: the monolith (Slice~1), a manually-engineered microservices
reference (Slice~2), the workflow's automated decomposition
(Slice~1~Refactored), and a Terraform configuration for Oracle Cloud
Infrastructure (Slice~3).

\textbf{What we measure.}
The paper reports empirical results from four measurement campaigns:
(i) structural alignment between the workflow's decomposition and the
independently-engineered manual reference, scored mechanically against
a pre-registered twenty-dimension rubric over ten repeated runs;
(ii) wall-clock duration of the automated decomposition compared to
the developer-time spent producing the manual reference;
(iii) an ablation that isolates the contribution of the skill scripts
above raw LLM capability by re-running the same model on the same
input with all skills disabled (ten baseline runs); and
(iv) precision and recall of the workflow's automated security
checks on a 20-file matched-pair injection corpus covering ten common
.NET vulnerability classes.

\textbf{Reproducibility.}
The skill scripts, the four artefact repositories tagged at the
experimental commit, the raw run logs and transcripts, the scoring
scripts, the verifier changelog, the security injection corpus, and
a full environment manifest (LLM model identity, CLI version,
hardware, dependency versions) are deposited at Zenodo and cited by
DOI in the manuscript.  The exact configuration can be reproduced
under the same model identity; transfer to other monoliths, languages
or models requires re-evaluation, and we discuss the limits of the
case-study scope explicitly in the threats-to-validity section.

\textbf{Scope and limits.}
We present the workflow as a reproducible pilot, not as a validated
general-purpose migration tool.  We make no claim that the decomposition
quality, productivity figures, or security-check performance generalise
beyond the case study without additional evaluation.  The contribution
is the composition pattern plus the reproducibility deposit, evaluated
honestly on one concrete .NET codebase.

\textbf{Fit for \textit{IEEE Access}.}
\textit{IEEE Access} accommodates reproducible case studies that
combine engineering practice with empirical evaluation, which is the
shape of this work.  The manuscript is structured around a clear
methodology section, multi-run empirical measurements with mean and
standard deviation, a baseline isolation experiment, an ablation
of failure modes, and a precision-recall security evaluation.

\textbf{Authors and contributions.}
\begin{itemize}
  \item \textbf{Florin Olariu} (Alexandru Ioan Cuza University of
    Ia\c{s}i and Suvoda LLC) --- skill design and implementation,
    experimental harness, evaluation, writing.
  \item \textbf{Lenu\c{t}a Alboaie} (Alexandru Ioan Cuza University of
    Ia\c{s}i) --- manual microservices reference (Slice~2) produced
    independently of the skill prompts, scoring rubric design,
    independent verification pass, writing.
\end{itemize}

\textbf{Originality and ethics.}
This work has not been published elsewhere and is not under review
at any other venue.  No human or animal subjects are involved.  We
declare no conflicts of interest.

We thank you for your consideration and look forward to the
reviewers' assessment.

\closing{Respectfully yours,}

\end{letter}
\end{document}
